% 调和函数（综述）
% license CCBYSA3
% type Wiki

本文根据 CC-BY-SA 协议转载翻译自维基百科\href{https://en.wikipedia.org/wiki/Harmonic_function}{相关文章}。


在数学、数理物理以及随机过程理论中，调和函数（harmonic function）是指满足拉普拉斯方程（Laplace's equation）的函数。  
设
\[
f : U \to \mathbb{R}~
\]
其中 \(U \subseteq \mathbb{R}^n\) 为一个开集，且 \(f\) 是定义在 \(U\) 上的二阶连续可微函数（即 \(f \in C^2(U)\)）。若 \(f\) 在 \(U\) 上处处满足
\[
\frac{\partial^2 f}{\partial x_1^2}
+ \frac{\partial^2 f}{\partial x_2^2}
+ \cdots
+ \frac{\partial^2 f}{\partial x_n^2}
= 0~
\]
则称 \(f\) 为一个\textbf{调和函数}。  
该方程通常记作
\[
\nabla^2 f = 0,
\quad \text{或等价地} \quad
\Delta f = 0.~
\]
\subsection{“Harmonic” 一词的词源}
“调和函数（harmonic function）”中的形容词“harmonic”（意为“谐和的”或“调和的”）起源于一根被拉紧的弦上某一点所经历的谐振运动（harmonic motion）。描述此类运动的微分方程的解可以用正弦函数（sine）与余弦函数（cosine）来表示，因此这些函数被称为谐波（harmonics）。傅里叶分析（Fourier analysis）研究如何将单位圆上的函数展开为这些谐波的级数。若考虑高维空间中单位 \(n\) 球面上的谐波类比，便得到所谓的球谐函数（spherical harmonics）。这些函数满足拉普拉斯方程。随着时间的推移，“harmonic”一词被推广，用来指代所有满足拉普拉斯方程的函数。\(^\text{[1]}\)  
\subsection{例子}
以下给出了一些二元调和函数的典型例子：
\begin{enumerate}
  \item \textbf{任意全纯函数（holomorphic function）的实部或虚部。}  
  实际上，定义在平面上的所有调和函数都可以表示为某个全纯函数的实部或虚部。
  \item 函数$f(x, y) = e^{x} \sin y$.这是前述情况的一个特例，因为$f(x, y) = \operatorname{Im}\!\left(e^{x + i y}\right)$,而 \( e^{x + i y} \) 是一个全纯函数。对 \(x\) 的二阶偏导数为$\frac{\partial^2 f}{\partial x^2} = e^{x} \sin y$,对 \(y\) 的二阶偏导数为$\frac{\partial^2 f}{\partial y^2} = -\, e^{x} \sin y$,因此$\frac{\partial^2 f}{\partial x^2} + \frac{\partial^2 f}{\partial y^2} = 0$,故 \(f\) 是调和的。
  \item 函数$f(x, y) = \ln\!\left(x^2 + y^2\right),\quad (x, y) \in \mathbb{R}^2 \setminus \{0\}$.该函数定义在去掉原点的平面上，可用于描述：一条线电荷（line charge）产生的电势；或一根无限长圆柱体质量（cylindrical mass）产生的引力势。
\end{enumerate}
对于三元调和函数，常见的例子列于下表。设$r^2 = x^2 + y^2 + z^2$,则可构造出许多满足拉普拉斯方程$\Delta f = 0$的函数，例如球对称势函数与球谐函数。

\begin{table}[ht]
\centering
\caption{}\label{tab_THhs1}
\begin{tabular}{|c|c|}
\hline
\textbf{函数}&\textbf{奇点}\\
\hline
\(\displaystyle \frac{1}{r}\)& 位于原点的单位点电荷 \\
\hline
\(\displaystyle \frac{x}{r^{3}}\)& 位于原点的沿 \(x\) 方向偶极子 \\
\hline
\(\displaystyle -\ln\!\left(r^{2}-z^{2}\right)\) & 整个 \(z\) 轴上的单位线电荷 \\
\hline
\(\displaystyle -\ln(r+z)\) & 负 \(z\) 轴上的单位线电荷 \\
\hline
\(\displaystyle \frac{x}{r^{2}-z^{2}}\) & 整个 \(z\) 轴上的沿 \(x\) 方向偶极子线 \\
\hline
\(\displaystyle \frac{x}{r(r+z)}\) & 负 \(z\) 轴上的沿 \(x\) 方向偶极子线 \\
\hline
\end{tabular}
\end{table}
在物理学中出现的调和函数（harmonic functions）由其奇点（singularities）和边界条件（boundary conditions）决定，例如狄利克雷边界条件（Dirichlet boundary conditions）或诺伊曼边界条件（Neumann boundary conditions）。在无边界区域上，如果给定一个调和函数，可以通过在其上加上任意整函数（entire function）的实部或虚部，得到具有相同奇点的另一个调和函数。因此，在这种情况下，调和函数并不由其奇点唯一确定；然而，在物理情形中，我们通常要求当\( r \to \infty \) 时，解趋于 \(0\)。在这种情况下，解的唯一性由刘维尔定理（Liouville's theorem）保证。

上述调和函数的奇点可以用静电学（electrostatics）的术语描述为“电荷”（charges）和“电荷密度”（charge densities），因此相应的调和函数与这些电荷分布产生的静电势（electrostatic potential）成比例。若将任意上述函数乘以常数、旋转，或加上常数，所得仍是一个调和函数。对这些函数作反演（inversion），将得到另一个调和函数，其奇点为原奇点关于某个球面“镜像”的位置。此外，任意两个调和函数的和仍然是一个调和函数。

最后，若考虑 \(n\) 个变量的情形，调和函数的例子包括：
\begin{enumerate}
\item 在整个 \(\mathbb{R}^{n}\) 上的常数函数、线性函数与仿射函数（affine functions），例如电容器极板间的电势，或平板质量分布产生的引力势；
\item 定义在 \(\mathbb{R}^{n}\setminus\{0\}\) 上的函数$f(x_{1}, \dots, x_{n}) = \left(x_{1}^{2} + \cdots + x_{n}^{2}\right)^{1 - n/2}, \quad n > 2$,它也是一个调和函数。
\end{enumerate}
\subsection{性质}
在给定开集 \( U \) 上的所有调和函数（harmonic functions）的集合，可以看作拉普拉斯算子 \(\Delta\) 的核，因此它构成一个实数域 \(\mathbb{R}\) 上的向量空间：任意调和函数的线性组合仍然是调和函数。

若 \( f \) 是定义在 \( U \) 上的调和函数，则 \( f \) 的所有偏导数在 \( U \) 上仍是调和函数。拉普拉斯算子 \(\Delta\) 与偏导算子在此类函数上可交换。

在许多方面，调和函数是解析函数（holomorphic functions）的实数对应物。所有调和函数都是解析的（analytic），即它们在局部可表示为幂级数展开。这一事实是椭圆算子（elliptic operators）的一般性质，而拉普拉斯算子是其中最典型的例子。

收敛的调和函数序列的逐点一致极限仍然是调和函数。这是因为任意满足平均值性质（mean value property）的连续函数都是调和的。例如，在区域 \((-\infty, 0) \times \mathbb{R}\) 上定义序列$f_{n}(x,y) = \frac{1}{n} e^{nx} \cos(ny)$,该序列中的每个函数都是调和的，并一致收敛于零函数。但注意其偏导数并不一致收敛于零函数（即零函数的导数）。此例说明在判断极限函数是否调和时，依赖平均值性质与连续性的重要性。
\subsection{与复变函数论的联系}
任意全纯函数（holomorphic function）的实部与虚部都是 \(\mathbb{R}^2\) 上的调和函数（称为一对共轭调和函数，harmonic conjugates）。反之，任意定义在 \(\mathbb{R}^2\) 的开子集 \(\Omega\) 上的调和函数 \(u\)，在局部都是某个全纯函数的实部。事实上，若写作 \( z = x + i y \)，则定义复函数\(g(z) := u_x - i u_y\),因为它满足柯西–黎曼方程（Cauchy–Riemann equations），所以 \( g \) 在 \(\Omega\) 内是全纯的。因此，\( g \) 在局部有一个原函数 \( f \)，并且 \( f' = g \)，从而 \( u \) 是 \( f \) 的实部（差一个常数）。

尽管上述与全纯函数的一一对应关系仅在两个实变量的情形成立，但 \( n \) 个变量的调和函数仍保留了许多典型的解析函数性质。例如：  它们是实解析的（real analytic）；满足最大值原理（maximum principle）与平均值原理（mean-value principle）；存在奇点可去定理（theorem of removal of singularities）与刘维尔定理（Liouville theorem），其形式与复变函数论中的相应定理完全类似。
\subsection{调和函数的性质}
许多调和函数的重要性质可以直接从拉普拉斯方程（Laplace's equation）推导出来。
\subsubsection{调和函数的正则性定理}
调和函数在开集上是无限可微的（infinitely differentiable）。实际上，调和函数是实解析函数（real analytic）的，也就是说，它们在局部可以表示为收敛幂级数。
\subsubsection{最大值原理}
调和函数满足以下最大值原理：若 \( K \subset U \) 是非空紧集（compact subset），则函数 \( f \) 在 \( K \) 上的限制 \( f|_{K} \) 的最大值与最小值都在 \( K \) 的边界上取得。若 \( U \) 是连通的（connected），则这意味着：调和函数 \( f \) 不可能在内部点处取得局部极大值或局部极小值，除非 \( f \) 是常数函数。类似的性质也可以推广到次调和函数（subharmonic functions）。
